\chapter{Conclusioni e Sviluppi Futuri}

Il presente capitolo ha la funzione di raccogliere le considerazioni conclusive relative al lavoro di ricerca e sviluppo svolto. L'elaborato si è mosso a partire da una fase di studio teorico sui protocolli di \textit{Zero Knowledge Proof}, passando in seguito alla loro sottoclasse dei protocolli Sigma, per poi concentrarsi in maniera specifica sul protocollo di identificazione di Schnorr. Di quest'ultimo sono stati analizzati sia il meccanismo di autenticazione interattiva, sia l'adattamento che consente di derivarne uno schema di firma digitale. Tale percorso ha fornito le basi concettuali necessarie per la progettazione e realizzazione dell'applicazione \textit{SchnorrAuthApp},
che integra questi principi in un contesto applicativo concreto.

Nonostante l'elevato livello di sicurezza garantito da un protocollo di autenticazione basato su meccanismi di \textit{Zero Knowledge Proof} (ZKP), la sua adozione in contesti reali risulta ancora complessa. I protocolli di autenticazione standard, come PAP o CHAP, continuano infatti a essere ampiamente utilizzati, poiché offrono un compromesso ritenuto accettabile tra sicurezza, semplicità di implementazione e facilità d'uso. Di conseguenza, i protocolli ZKP, pur rappresentando una soluzione tecnologicamente più avanzata, faticano a trovare un impiego diffuso al di fuori di scenari sperimentali o ambiti in cui la sicurezza riveste un'importanza critica.

L'applicazione sviluppata rappresenta comunque un risultato significativo, poiché dimostra la fattibilità tecnica di un sistema di autenticazione fondato su principi crittografici avanzati. \textit{SchnorrAuthApp} traduce in termini concreti e operativi un modello teorico complesso, mostrando come i protocolli ZKP possano essere implementati in modo efficace anche in ambienti mobili e multipiattaforma, grazie al framework Flutter. La pubblicazione del codice sorgente su una piattaforma open source come GitLab contribuisce inoltre alla trasparenza e alla riproducibilità del lavoro, favorendo la collaborazione e la possibilità di ulteriori evoluzioni da parte della comunità.

Durante lo sviluppo sono emerse tuttavia alcune sfide tecniche e concettuali. Una delle principali difficoltà è legata alla gestione sicura delle chiavi crittografiche, in particolare nei contesti mobili dove la persistenza dei dati e la protezione dell'ambiente di esecuzione non sono sempre garantite. Allo stesso modo, l'esigenza di bilanciare sicurezza e usabilità rappresenta un aspetto centrale: un sistema di autenticazione, per essere realmente adottato, deve infatti risultare semplice da usare e da integrare con le infrastrutture esistenti. Inoltre, la mancanza di standard consolidati per l'integrazione dei protocolli ZKP nei framework di autenticazione più diffusi (come OAuth 2.0 o OpenID Connect) costituisce un ulteriore ostacolo alla loro diffusione su larga scala.

\section{Possibili Sviluppi Futuri}

Gli sviluppi futuri del progetto potranno concentrarsi su diverse direzioni. Un primo passo potrà consistere nell'integrazione del protocollo di Schnorr con architetture di autenticazione standard, in modo da ampliare la compatibilità del sistema con applicazioni web e servizi distribuiti.
Un'evoluzione naturale del progetto potrebbe consistere nell'integrazione del sistema di autenticazione basato sul protocollo di Schnorr all'interno di un server web, come componente o plugin di librerie e framework HTTP comunemente utilizzati nello sviluppo di applicazioni web. In questo modo, il protocollo potrebbe essere impiegato non solo per l'autenticazione degli utenti, ma anche per la protezione delle API, offrendo un'alternativa sicura ai tradizionali meccanismi basati su token o password. Una tale integrazione permetterebbe di estendere l'utilizzo del sistema a contesti reali e distribuiti, sfruttando l'interoperabilità offerta dai principali linguaggi di programmazione e framework moderni.

Un'ulteriore estensione potrebbe riguardare l'introduzione di una funzionalità per l'esportazione e l'importazione sicura della chiave privata, utile a creare copie di backup protette.
In questo modo, un utente che dispone di un solo dispositivo avrebbe la possibilità di salvare la propria chiave segreta su un supporto esterno, come un hard drive, senza comprometterne la sicurezza.
Tale meccanismo garantirebbe una maggiore resilienza ai guasti o alla perdita del dispositivo principale, mantenendo al contempo la riservatezza della chiave segreta.
